\documentclass[10pt,a4paper]{amsart}
\usepackage[margin=19mm]{geometry}
\usepackage{amsmath,amssymb,mathtools}
\usepackage[scr=rsfs,cal=euler]{mathalfa}
\usepackage{xfrac}
\usepackage[unicode,hidelinks,bookmarks=false]{hyperref}
\newcommand{\ie}{i.e.\ }
\newcommand{\cf}{cf.\ }
\newcommand{\resp}{respectively\ }
\newcommand{\Subsection}{\subsection}
\providecommand{\nobreakdash}{\nobreak\hbox{-}}
\setlength{\emergencystretch}{2em}
\multlinegap0pt
\usepackage{mathtools}
\usepackage{dsfont}
\usepackage{amssymb}
\newtheorem{prop}{Proposition}
\newtheorem{lem}{Lemma}
\def\C{\mathbb{C}}
\newcommand{\End}{\operatorname{End}}
\newcommand{\Hom}{\operatorname{Hom}}
\newcommand{\KL}{KL}
\newcommand{\id}{\operatorname{id}}
\newcommand{\rad}{\operatorname{rad}}
\title{Exactness and Drinfeld--Sokolov Realization of the McRae--Yang Tensor Functors}
\author{Shun Xu}
\date{Source: arXiv:2608.13869v2 (CC BY 4.0)}
\begin{document}
\maketitle

\noindent\emph{Source and licence.} Exactness and Drinfeld--Sokolov Realization of the McRae--Yang Tensor Functors. Version arXiv:2608.13869v2 (CC BY 4.0), distributed by arXiv under the Creative Commons Attribution 4.0 International licence, http://creativecommons.org/licenses/by/4.0/, as recorded in the arXiv licence metadata for this exact version and shown on the abstract page https://arxiv.org/abs/2608.13869v2. Full licence notice: \texttt{licenses/arxiv-2608.13869-CC-BY-4.0.txt}.\par\smallskip

\noindent\textit{Editorial scope.} Counted unit: the article's prop prop:natural-on-projectives (source0.tex lines 6316--6325). The article's complete original proof of it is delivered with it (source0.tex lines 6327--6438, one \texttt{proof} environment of 112 lines). No mathematics is written, repaired or re-derived: the statement and the proof are the article's own lines, in the article's own notation, and no retained line is substituted or re-flowed.  Retained uncounted as the source-local context the proof needs: lines 331--335; lines 362--393; lines 6214--6226; lines 6259--6273; lines 6290--6296.  Nothing was omitted from any retained span, so the package carries no omission record.  No retained line contains a citation, so the wrapper carries no bibliography and no bibliography entry.  No retained line contains a figure environment, an image-inclusion command or drawing code, so no image is delivered and the package declares no asset dependency.  Editorial formatting only, in addition: an AMS \texttt{amsart} preamble that restates the article's own declarations and macros used by the retained lines, namely the environments \texttt{prop}, \texttt{lem} on separate counters, together with the standard packages those lines need.  The retained environments print in this document as Proposition 1, Lemma 1 in order of appearance, whereas the article prints the same items with the numbering of its own full document.  The complete native source of the article as distributed in the arXiv e-print for this exact version is bundled unchanged as \texttt{sources/207618/source0.tex}.
\begin{equation}\label{eq:reflection-chain}
 s_{2m}(r)=2mp+r,\qquad
 s_{2m+1}(r)=2(m+1)p-r
 \qquad(m\geq0).
\end{equation}

\begin{prop}[Source projective Hom package]\label{prop:source-Hom-package}
For a fixed reflection chain \eqref{eq:reflection-chain},
\[
 \dim\Hom(P_{s_i},P_{s_j})=
 \begin{cases}
 2,&i=j\geq1,\\
 1,&i=j=0,\\
 1,&|i-j|=1,\\
 0,&|i-j|\geq2.
 \end{cases}
\]
If $P_{s_i(r)}$ and $P_{s_j(r')}$ belong to distinct reflection chains,
then both Hom spaces between them vanish.  The wall projectives $P_{np}$ are isolated
simple-projectives:
\[
 \End(P_{np})=\C,\qquad
 \Hom(P_{np},P_j)=\Hom(P_j,P_{np})=0\quad(j\ne np).
\]
In particular, $\Hom(P_{np},P_{mp})=0$ for $n\ne m$.

Moreover, for each adjacent pair one may choose nonzero arrows
\[
 x_i:P_{s_i}\longrightarrow P_{s_{i+1}},
 \qquad
 y_i:P_{s_{i+1}}\longrightarrow P_{s_i}
\]
such that
\begin{equation}\label{eq:source-zigzag-composition}
 \nu_{i+1}:=x_i\circ y_i
\end{equation}
is the nonzero radical endomorphism of $P_{s_{i+1}}$.
\end{prop}

\begin{lem}[Intrinsic conformal-radical compatibility]
\label{lem:twist-locking}
Let $P=P_{np+r}$ with $n\ge1$ and $1\le r\le p-1$.  Every
Virasoro-module isomorphism
\[
 \eta_P:F(P)\xrightarrow{\sim}H(P)
\]
satisfies
\begin{equation}\label{eq:twist-locking}
 \eta_PF(u)=H(u)\eta_P
 \qquad\text{for every }u\in\rad\End(P).
\end{equation}
\end{lem}

\begin{lem}[Mixed projective Hom dimensions]
\label{lem:mixed-projective-Hom}
Fix a non-wall reflection chain $\{s_i\}_{i\ge0}$ and put
$\Pi_i=P_{s_i}$.  For any $\mathscr G,\mathscr G'\in\{F,H\}$,
\[
 \dim\Hom(\mathscr G(\Pi_i),\mathscr G'(\Pi_j))
 =
 \begin{cases}
  2,&i=j\ge1,\\
  1,&i=j=0,\\
  1,&|i-j|=1,\\
  0,&|i-j|\ge2.
 \end{cases}
\]
\end{lem}

\begin{lem}[Projective exhaustion]\label{lem:projective-exhaustion}
Every indecomposable projective object of $\KL^k(\mathfrak{sl}_2)$ is
isomorphic to exactly one $P_j$.  Consequently
\[
 \mathcal P^k=\operatorname{add}\{P_j:j\ge1\}.
\]
\end{lem}

\begin{prop}\label{prop:natural-on-projectives}
There is a natural isomorphism of $\C$-linear functors
\[
 \eta:F|_{\mathcal P^k}\xRightarrow{\ \sim\ }H|_{\mathcal P^k}.
\]
On each non-wall reflection chain the components can be chosen
inductively.  Exponential compatibility forces the scalar obstruction for
the reverse adjacent morphisms to be trivial; no uniqueness of the
individual components is asserted.
\end{prop}

\begin{proof}
Fix a non-wall reflection chain $\{s_i\}_{i\ge0}$, put
$\Pi_i=P_{s_i}$, and choose adjacent generators
\[
 \Pi_i\underset{y_i}{\stackrel{x_i}{\rightleftarrows}}\Pi_{i+1}
\]
as nonzero generators of the two one-dimensional adjacent Hom spaces in
Proposition~\ref{prop:source-Hom-package}, normalized only so that
\begin{equation}\label{eq:zigzag-loop-section6}
 \nu_{i+1}:=x_i\circ y_i\ne0
\end{equation}
spans $\rad\End(\Pi_{i+1})$.

Choose an isomorphism $\eta_0:F(\Pi_0)\to H(\Pi_0)$.  Suppose $\eta_i$
has been chosen and take any isomorphism
$\eta'_{i+1}:F(\Pi_{i+1})\to H(\Pi_{i+1})$.  By
Lemma~\ref{lem:mixed-projective-Hom}, both
$\eta'_{i+1}\circ F(x_i)$ and $H(x_i)\circ\eta_i$ lie in the same one-dimensional
mixed Hom space.  They are nonzero: $F(x_i)\ne0$ by projective faithfulness of $F$,
$H(x_i)\ne0$ by faithfulness of $H$, and both $\eta_i$ and
$\eta'_{i+1}$ are isomorphisms.  Hence there is a unique
$a_i\in\C^\times$ such that
\[
 \eta'_{i+1}\circ F(x_i)=a_i H(x_i)\circ\eta_i.
\]
Set $\eta_{i+1}:=a_i^{-1}\eta'_{i+1}$.  Then
\begin{equation}\label{eq:forward-naturality}
 \eta_{i+1}\circ F(x_i)=H(x_i)\circ\eta_i.
\end{equation}
The reverse mixed Hom space is also one-dimensional.  Both
$\eta_i\circ F(y_i)$ and $H(y_i)\circ\eta_{i+1}$ are nonzero, by
faithfulness of $F$ on projectives, faithfulness of $H$, and invertibility
of the chosen components.  Hence
\begin{equation}\label{eq:reverse-scalar}
 \eta_i\circ F(y_i)=c_iH(y_i)\circ\eta_{i+1}
\end{equation}
for a unique $c_i\in\C^\times$.  Composing with the forward relation gives
\begin{align*}
 \eta_{i+1}\circ F(\nu_{i+1})
 &=\eta_{i+1}\circ F(x_i)\circ F(y_i)\\
 &=H(x_i)\circ\eta_i\circ F(y_i)\\
 &=c_iH(x_i)\circ H(y_i)\circ\eta_{i+1}\\
 &=c_iH(\nu_{i+1})\circ\eta_{i+1}.
\end{align*}
Lemma~\ref{lem:twist-locking} gives
$\eta_{i+1}\circ F(\nu_{i+1})=
H(\nu_{i+1})\circ\eta_{i+1}$.  Subtracting from the
preceding equality gives
\[
 (c_i-1)H(\nu_{i+1})\circ\eta_{i+1}=0.
\]
Since $\nu_{i+1}\ne0$, faithfulness of $H$ gives
$H(\nu_{i+1})\ne0$, and hence
$H(\nu_{i+1})\circ\eta_{i+1}\ne0$.  Therefore $c_i=1$.  Thus the normalization along $x_i$ automatically
enforces naturality along the reverse arrow $y_i$ as well.  This argument
uses only the fixed nonzero loop $\nu_{i+1}=x_i y_i$; no assertion about
the opposite composite $y_i x_i$ is required.

By Proposition~\ref{prop:source-Hom-package}, the displayed elements form
bases of every nonzero Hom space between indecomposable projectives on the
fixed reflection chain:
\[
 \C\id_{\Pi_0},\qquad \C x_i,\qquad \C y_i,
 \qquad \C\id_{\Pi_i}\oplus\C\nu_i\ (i\ge1),
\]
and all remaining Hom spaces vanish.  Naturality has been checked on these
bases: for $x_i,y_i$ above, for $\nu_i$ by
Lemma~\ref{lem:twist-locking}, and for identities automatically.  Linearity
therefore gives naturality for every morphism on the chain.  Hom spaces between distinct reflection chains vanish by
Proposition~\ref{prop:source-Hom-package}.

For every wall projective $P_{np}$, choose any isomorphism
\[
 \eta_{np}:F(P_{np})\xrightarrow{\sim}H(P_{np}),
\]
using the identifications
$F(P_{np})\cong K_{np}\cong H(P_{np})$.  Wall projectives are isolated and
have scalar endomorphism algebras, so naturality there is automatic.
By Lemma~\ref{lem:projective-exhaustion}, these components cover all
indecomposable projective objects.

For $R=0$ take the unique isomorphism $0\to0$.  For each
indecomposable $P_j$ set $\alpha_{P_j}=\id_{P_j}$.  For every other
nonzero projective $R$, fix once and for all an isomorphism
\[
 \alpha_R:R\xrightarrow{\sim}
 \bigoplus_j P_j^{\oplus m_j(R)}
\]
with only finitely many nonzero multiplicities.  No independence of this
choice is asserted or needed.  Since $F$ and $H$ are additive, we use their canonical finite-biproduct identifications
\[
 F\!\left(\bigoplus_jP_j^{\oplus m_j(R)}\right)
 \cong\bigoplus_jF(P_j)^{\oplus m_j(R)},
 \qquad
 H\!\left(\bigoplus_jP_j^{\oplus m_j(R)}\right)
 \cong\bigoplus_jH(P_j)^{\oplus m_j(R)}.
\]
With these identifications understood, define
\[
 \eta_R:=
 H(\alpha_R)^{-1}
 \left(\bigoplus_j\eta_{P_j}^{\oplus m_j(R)}\right)
 F(\alpha_R).
\]
No canonicity or independence of the chosen Krull--Schmidt isomorphisms is
asserted.  Fix these choices.  We now verify that the resulting family is
natural.  If $f:R\to R'$ is any projective morphism, the matrix of
$\alpha_{R'}f\alpha_R^{-1}$ has entries between indecomposable
projectives, on which naturality has already been proved.  Matrix
multiplication therefore gives
$\eta_{R'}F(f)=H(f)\eta_R$.
\end{proof}

\end{document}
